\documentclass[11pt,a4paper]{article}
\usepackage[english]{cursus}
\author{dr. Maarten Slembrouck}

\begin{document}

\labtitle{Lab 4 -- Data Processing II}{Combine, summarize, reshape}{Course text for the notebook \code{lab4\_data\_processing\_II.py}}

\section*{What is this lab about?}

Last week you turned messy exports into tidy tables: one row per
athlete, one row per test. But a coach's questions are rarely about one
row. ``How tall are our defenders on average?'' is about a \emph{group}
of rows. ``How fast do our goalkeepers sprint?'' needs information from
\emph{two} tables. ``Has everyone improved since the first test?'' asks
for the table to be \emph{arranged} differently. And the heart-rate data
of a training session is a \emph{time series}, with one measurement per
second.

\textbf{After this lab you can} summarize data per group with
\code{groupby}, combine tables with \code{merge}, reshape data with
\code{melt} and \code{pivot\_table}, and compute basic features of a
time series (heart-rate zones, rolling average, resampling).

Just like last week, every example was actually run on the club's data.
Read this text next to the notebook and type the examples yourself.

% ---------------------------------------------------------------------------
\section{Starting point: last week's cleaned data}

The notebook starts with the cleaning from lab 3, in compact form. Two
helper functions do the work: \code{clean\_positions} turns the 13
spellings into 4 positions, \code{fix\_sprint\_units} converts
milliseconds into seconds.

\begin{python}[noout]
import pandas as pd

def clean_positions(position_column):
    position_map = {
        "forward": "Forward", "fwd": "Forward",
        "midfielder": "Midfielder", "mid": "Midfielder",
        "defender": "Defender", "def": "Defender",
        "goalkeeper": "Goalkeeper", "gk": "Goalkeeper",
    }
    return position_column.str.strip().str.lower().map(position_map)

def fix_sprint_units(seconds_column):
    return seconds_column.where(seconds_column < 20, seconds_column / 1000)

REFERENCE_DATE = pd.Timestamp("2026-09-14")

athletes = pd.read_csv(f"{DATA_DIR}/athletes.csv").assign(
    position=lambda d: clean_positions(d["position"]),
    weight_kg=lambda d: d["weight_kg"].str.replace(",", ".").astype(float),
    age=lambda d: (REFERENCE_DATE - pd.to_datetime(d["birthdate"])).dt.days // 365,
)
tests = pd.read_csv(f"{DATA_DIR}/tests.csv").assign(
    sprint_10m_s=lambda d: fix_sprint_units(d["sprint_10m_s"]),
    sprint_30m_s=lambda d: fix_sprint_units(d["sprint_30m_s"]),
)
\end{python}

Read this code calmly: you saw everything in it last week. Note that the
duplicated athlete is \emph{deliberately} not removed yet. That becomes
important in section~\ref{sec:merge}.

\subsection*{Recap: reading error messages}

The notebook starts with a broken cell that you have to repair. An error
message looks scary, but the most important information is always on
the \textbf{last line}: the kind of error and a short explanation. Two
errors you will certainly meet today:

\begin{python}[error]
athletes["positio"]
\end{python}

A \code{KeyError} means: this key (here: column name) does not exist.
Almost always a typo.

\begin{python}[error]
pd.Series(["57,5", "60,3"]).astype(float)
\end{python}

A \code{ValueError} from \code{astype(float)} means: I can't turn this
text into a number. Here the comma is in the way.

% ---------------------------------------------------------------------------
\section{Summarizing groups with \code{groupby}}

\subsection{Split, apply, combine}

Suppose the coach wants the average age per position. In your head you
do three things:

\begin{enumerate}
  \item \textbf{split}: divide the athletes into four groups, one per
        position;
  \item \textbf{apply}: compute the mean for each group;
  \item \textbf{combine}: put the four means in one table.
\end{enumerate}

\begin{center}
\begin{tikzpicture}[font=\small\sffamily, box/.style={rounded corners=2pt, minimum width=2.3cm, minimum height=0.6cm, text=white}]
  \node[draw=none, fill=tint, rounded corners=3pt, minimum width=2.3cm, minimum height=3.2cm, text=pitch, align=center] (all) at (0,0) {all\\athletes};
  \foreach \p/\c/\y in {Forward/forward/1.2, Midfielder/midfielder/0.4, Defender/defender/-0.4, Goalkeeper/goalkeeper/-1.2} {
    \node[box, fill=\c] (\p) at (3.6,\y) {\p};
    \draw[-{Stealth}, muted] (all.east) -- (\p.west);
    \node[text=muted, font=\small\itshape] (m\p) at (6.1,\y) {mean};
    \draw[-{Stealth}, muted] (\p.east) -- (m\p.west);
    \draw[-{Stealth}, muted] (m\p.east) -- (8.2,0);
  }
  \node[draw=none, fill=pitch, rounded corners=3pt, minimum width=2.1cm, minimum height=2.0cm, text=white, align=center] at (9.3,0) {one row\\per position};
  \node[text=pitch, font=\small\bfseries\sffamily] at (1.8,2.1) {split};
  \node[text=pitch, font=\small\bfseries\sffamily] at (5.0,2.1) {apply};
  \node[text=pitch, font=\small\bfseries\sffamily] at (8.4,2.1) {combine};
\end{tikzpicture}
\end{center}

In pandas that is one line. \code{groupby("position")} splits,
\code{["age"]} picks the column, \code{.mean()} applies the mean and
combines:

\begin{python}
mean_age_per_position = athletes.groupby("position")["age"].mean().round(1)
mean_age_per_position
\end{python}

The result is a Series with the positions as its \emph{index}. Instead
of \code{.mean()} you can also use \code{.median()}, \code{.min()},
\code{.max()}, \code{.sum()} or \code{.count()}.

\begin{watchout}[\code{size()} or \code{count()}?]
\code{.size()} counts the rows per group. \code{.count()} counts the
\emph{filled-in} values in a column and skips \code{NaN}. To know how
many athletes there are per position, use \code{.size()}.
\end{watchout}

\begin{python}
athletes.groupby("position").size()
\end{python}

\subsection{Numbers into categories: \code{pd.cut}}

Youth teams are divided into age groups: U13, U15, U17, U19.
\code{pd.cut} divides a numeric column into bins. You give the edges
(\code{bins}) and the names of the bins (\code{labels}):

\begin{python}
athletes_with_age_group = athletes.assign(
    age_group=lambda d: pd.cut(
        d["age"], bins=[10, 12, 14, 16, 18], labels=["U13", "U15", "U17", "U19"]
    )
)
athletes_with_age_group[["name", "age", "age_group"]].head()
\end{python}

Five edges make four bins: from 10 to 12, from 12 to 14, and so on. A
bin includes its \emph{right} edge but not its left edge: an athlete of
12 falls in U13 (10 up to and including 12), an athlete of 13 in U15. So
there is always one label fewer than there are edges.

\subsection{Grouping on several columns and \code{agg}}

You can group on two columns at once by passing a list. With
\code{.agg(...)} you compute several statistics in one go and give each
new column a name yourself:

\begin{python}
athletes_with_age_group.groupby(["position", "age_group"], observed=True).agg(
    n_athletes=("athlete_id", "size"),
    mean_height_cm=("height_cm", "mean"),
)
\end{python}

Read \code{mean\_height\_cm=("height\_cm", "mean")} as: ``make a column
\code{mean\_height\_cm} by applying the function \code{mean} to the
column \code{height\_cm}''. A few extra things:

\begin{itemize}
  \item The index now has two levels (position and age group): a
        \emph{MultiIndex}. \code{.reset\_index()} turns them back into
        ordinary columns.
  \item \code{age\_group} is a \emph{categorical} column (made by
        \code{pd.cut}). \code{observed=True} makes sure combinations
        without athletes (e.g.\ goalkeepers in U13) don't show up as an
        empty row.
\end{itemize}

\begin{tryit}
Compute the maximum height per position. Then group by gender and count
with \code{.size()} how many boys and girls there are.
\end{tryit}

% ---------------------------------------------------------------------------
\section{Combining tables with \code{merge}}\label{sec:merge}

\subsection{A key in two tables}

\code{athletes} knows the position of every athlete, \code{tests} knows
the sprint times. To look at sprint time per position, you have to
combine the two. That works because both tables have an
\code{athlete\_id} column: the \textbf{key}.

A small example makes clear what happens:

\begin{python}
people = pd.DataFrame({"athlete_id": ["A1", "A2", "A3"],
                       "name": ["Yana", "Stan", "Emma"]})
results = pd.DataFrame({"athlete_id": ["A1", "A1", "A2", "A9"],
                        "sprint_s": [5.3, 5.2, 5.9, 6.1]})
people.merge(results, on="athlete_id", how="inner")
\end{python}

For every row in \code{people}, pandas looks for the rows in
\code{results} with the same \code{athlete\_id}. Yana has two results,
so she appears twice. Emma (A3) has no result and A9 has no name: with
\code{how="inner"} both disappear.

\begin{python}
people.merge(results, on="athlete_id", how="left")
\end{python}

With \code{how="left"}, \emph{every} row of the left table is kept. Emma
is there, with \code{NaN} as her sprint time.

\begin{keyidea}[The kinds of joins]
\textbf{inner}: only keys that are in both tables.
\textbf{left}: all rows of the left table, filled in where possible.
\textbf{right}: the same, from the right table.
\textbf{outer}: everything from both tables.
Check the number of rows after every merge: it tells you whether
something disappeared or doubled.
\end{keyidea}

With the real data:

\begin{python}
merged_inner = athletes.merge(tests, on="athlete_id", how="inner")
merged_left = athletes.merge(tests, on="athlete_id", how="left")
(len(athletes), len(tests), len(merged_inner), len(merged_left))
\end{python}

Here both have the same length, because every athlete has at least one
test result. With real data that is rarely guaranteed.

\subsection{The duplicated athlete}

One athlete appears twice in \code{athletes}, with two different IDs. You
find those rows with \code{duplicated}; with \code{keep=False}
\emph{all} copies are marked, not only the second one:

\begin{python}
same_person = athletes.duplicated(subset=["name", "birthdate"], keep=False)
athletes.loc[same_person, ["athlete_id", "name", "birthdate"]]
\end{python}

The merge does exactly what it was asked, but does not know these are
the same person. After the merge, that athlete's tests appear under two
IDs:

\begin{python}
dupe_name = athletes.loc[same_person, "name"].iloc[0]
merged_inner.loc[merged_inner["name"] == dupe_name,
                 ["athlete_id", "name", "test_moment", "sprint_30m_s"]]
\end{python}

Every statistic you compute after this merge counts this athlete twice.
The solution: remove the duplicates \emph{before} the merge, by name and
birthdate. After the merge that no longer works, because then the
different test moments of one athlete also share the same name and
birthdate: you would throw away real test results.

\begin{python}
athletes_unique = athletes.drop_duplicates(subset=["name", "birthdate"], keep="first")
combined = athletes_unique.merge(tests, on="athlete_id", how="inner")
combined.shape
\end{python}

\subsection{Sets and \code{isin}}

To find out which rows disappeared when removing the duplicate, you need
two tools.

A \textbf{set} is a collection without duplicates and without order. The
difference of two sets (\code{a - b}) gives the elements that are in
\code{a} but not in \code{b}:

\begin{python}
set(["A1", "A2", "A3"]) - set(["A2"])
\end{python}

\code{.isin(...)} gives a mask that is \code{True} for every row whose
value is in the given list or set:

\begin{python}
results["athlete_id"].isin({"A1", "A9"})
\end{python}

With these two building blocks (and \code{.sum()} on a mask) you can
solve exercise 2 in the notebook.

% ---------------------------------------------------------------------------
\section{Long and wide: \code{pivot\_table} and \code{melt}}

\subsection{Two shapes for the same data}

\code{combined} is \textbf{long}: one row per athlete \emph{per test
moment}. An athlete with three tests appears three times. The same
information can also be \textbf{wide}: one row per athlete, with a
column per test moment.

\begin{python}
sprint_wide = combined.pivot_table(
    index="athlete_id", columns="test_moment", values="sprint_30m_s"
)
sprint_wide.head()
\end{python}

Read the three arguments like this:
\begin{itemize}
  \item \code{index}: what becomes a row (one per athlete);
  \item \code{columns}: which values become columns (test moments 1, 2
        and 3);
  \item \code{values}: which value goes in the cells (the sprint time).
\end{itemize}

If a combination is missing, you get \code{NaN}. If there are several
rows for the same combination, \code{pivot\_table} takes the mean by
default.

\subsection{Why wide is handy}

In the wide shape, progress between two test moments is a single
subtraction. Note: the column names here are \emph{numbers}, not text.
So you write \code{sprint\_wide[2]}, not \code{sprint\_wide["2"]}.

\begin{python}
sprint_progress_1_to_2 = sprint_wide[2] - sprint_wide[1]
sprint_progress_1_to_2.describe()
\end{python}

A negative number means progress here: a lower sprint time is faster.
For jump height it is the other way round: higher is better.

\subsection{Back to long: \code{melt}}

\code{melt} does the opposite of \code{pivot\_table}: columns become rows
again. Plotly express (lab 5) prefers long data. First,
\code{.reset\_index()} turns the index back into an ordinary column:

\begin{python}
sprint_long_again = sprint_wide.reset_index().melt(
    id_vars="athlete_id", var_name="test_moment", value_name="sprint_30m_s"
)
sprint_long_again.head()
\end{python}

\code{id\_vars} are the columns that stay, \code{var\_name} is the name
of the new column holding the former column names, \code{value\_name}
the name of the column holding the values.

% ---------------------------------------------------------------------------
\section{Time series: heart rate during a training session}

\code{hr\_session.csv} contains the heart rate and speed of five athletes
per second during one training session. A table in which every row is a
moment in time is called a \textbf{time series}.

\subsection{Making a timestamp}

The file has a date and the number of seconds since the start
(\code{time\_s}). Together they make a timestamp.
\code{pd.to\_timedelta(..., unit="s")} turns a number of seconds into a
duration, which you can add to a date:

\begin{python}
hr_session = pd.read_csv(f"{DATA_DIR}/hr_session.csv")
hr_session["timestamp"] = pd.to_datetime(hr_session["date"]) + pd.to_timedelta(
    hr_session["time_s"], unit="s"
)
hr_session.head(3)
\end{python}

\subsection{Heart-rate zones}

Coaches express intensity as a percentage of the \emph{maximum} heart
rate. A common rule of thumb estimates it as 220 minus the age. The age
is in \code{athletes}, so first a merge:

\begin{python}
ages = athletes[["athlete_id", "age"]]
hr_with_age = hr_session.merge(ages, on="athlete_id", how="left")
hr_with_age["hr_max"] = 220 - hr_with_age["age"]
hr_with_age["pct_hr_max"] = (
    100 * hr_with_age["heart_rate_bpm"] / hr_with_age["hr_max"]
)
hr_with_age[["athlete_id", "heart_rate_bpm", "hr_max", "pct_hr_max"]].head(3)
\end{python}

\code{pd.cut} divides the percentage into five zones. The odd number
1000 as the last edge makes sure a percentage above 100 still falls in
zone 5:

\begin{python}
hr_zoned = hr_with_age.assign(
    hr_zone=lambda d: pd.cut(
        d["pct_hr_max"],
        bins=[0, 60, 70, 80, 90, 1000],
        labels=["Zone 1", "Zone 2", "Zone 3", "Zone 4", "Zone 5"],
    )
)
hr_zoned["hr_zone"].value_counts()
\end{python}

\subsection{Time per zone}

Every row is one second. So the number of rows per athlete and per zone,
divided by 60, gives the number of minutes:

\begin{python}
time_per_zone_minutes = (
    hr_zoned.groupby(["athlete_id", "hr_zone"], observed=True).size() / 60
).rename("minutes").reset_index()
time_per_zone_minutes.head(5)
\end{python}

If you only want certain zones, first filter the rows with
\code{.isin([...])} and then group. That is the core of exercise 4.

\subsection{Rolling average and resampling}

A heart rate per second jumps up and down. Two techniques make the
series calmer. Both need a time series \emph{sorted} by time, with the
time as the index (\code{set\_index("timestamp")}):

\begin{python}[noout]
one_athlete = hr_zoned.loc[hr_zoned["athlete_id"] == "A001"]
one_athlete = one_athlete.sort_values("timestamp")
hr_series = one_athlete.set_index("timestamp")["heart_rate_bpm"]
\end{python}

A \textbf{rolling average} (\code{rolling}) computes, at every moment,
the mean of the past 30 seconds. You keep as many rows, but the line
gets smoother:

\begin{python}
hr_series.rolling("30s").mean().iloc[[0, 1, 29, 30, 31]]
\end{python}

\textbf{Resampling} (\code{resample}) groups the time series into fixed
blocks, here of one minute, and summarizes each block. You get far fewer
rows: one per minute.

\begin{python}
hr_series.resample("1min").mean().head(4)
\end{python}

\begin{keyidea}[\code{rolling} or \code{resample}?]
\code{rolling} gives one value per original row (a sliding window);
useful to make a chart smoother. \code{resample} gives one value per
time block; useful to summarize per minute or per quarter of an hour.
You can think of \code{resample} as a \code{groupby} on time.
\end{keyidea}

% ---------------------------------------------------------------------------
\section{The mini-case: a training report}

In the mini-case (in pairs) you build one row per athlete with the
minutes in zone 4--5, the maximum speed and the distance covered. Two
ideas help you.

\textbf{Distance from speed.} Speed is in km/h and there is one
measurement per second. So in one second you cover
\code{speed\_kmh / 3600} kilometres. The sum of those distances over the
whole session is the total distance.

\textbf{Combining Series via the index.} When two Series have the same
index (e.g.\ the \code{athlete\_id}), pandas puts the values next to the
right row in an assignment or calculation, even if the order is
different. This is called \emph{aligning on the index}:

\begin{python}
report = pd.DataFrame({"max_speed": [27.1, 24.5]}, index=["A002", "A001"])
extra = pd.Series([10.0, 12.5], index=["A001", "A002"])
report["minutes"] = extra
report
\end{python}

A001 gets 10.0 and A002 gets 12.5, even though they were in a different
order in \code{extra} than in \code{report}.

% ---------------------------------------------------------------------------
\section{Summary}

\begin{center}
\rowcolors{2}{tint}{white}
\begin{tabularx}{\textwidth}{@{}L{8.4cm}Y@{}}
\toprule
\textbf{Code} & \textbf{What it does} \\
\midrule
\code{df.groupby("a")["b"].mean()} & mean of \code{b} per value of \code{a} \\
\code{df.groupby("a").size()} & number of rows per group \\
\code{df.groupby(["a", "b"]).agg(name=("col", "mean"))} & several named statistics per group \\
\code{pd.cut(col, bins=..., labels=...)} & divide numbers into categories \\
\code{.reset\_index()} & turn the index back into ordinary columns \\
\code{left.merge(right, on="key", how="inner")} & combine two tables on a key \\
\code{set(a) - set(b)}, \code{col.isin(list)} & difference of sets, ``is in'' mask \\
\code{df.pivot\_table(index=, columns=, values=)} & from long to wide \\
\code{df.melt(id\_vars=, var\_name=, value\_name=)} & from wide to long \\
\code{pd.to\_timedelta(col, unit="s")} & number of seconds to a duration \\
\code{s.rolling("30s").mean()} & rolling average over 30 seconds \\
\code{s.resample("1min").mean()} & mean per block of one minute \\
\bottomrule
\end{tabularx}
\end{center}

\subsection*{Common mistakes}

\begin{itemize}
  \item Not checking the number of rows after a merge.
  \item Removing duplicates \emph{after} the merge instead of before.
  \item Using \code{.count()} to count rows when there are missing
        values; use \code{.size()}.
  \item Writing column names as text after a \code{pivot\_table}
        (\code{wide["2"]}) when they are numbers (\code{wide[2]}).
  \item Using \code{rolling} or \code{resample} on data that is not
        sorted by time, or without a time index.
\end{itemize}

\subsection*{Glossary}

\begin{description}[leftmargin=3.2cm, labelwidth=3.0cm, font=\normalfont\bfseries]
  \item[groupby] split into groups, compute something per group, and combine the results.
  \item[aggregation] summarizing a group of values into one value (mean, sum, count).
  \item[MultiIndex] index with several levels, e.g.\ position and age group.
  \item[key] column used to combine two tables.
  \item[join] way of combining: inner, left, right or outer.
  \item[long / wide] one row per measurement, or one row per athlete with a column per measurement.
  \item[time series] table in which every row is a moment in time.
\end{description}

\end{document}
